\section*{Bab \ref{ch:g10:numbers} --- Bilangan dan Himpunan Bilangan}

\begin{solution}{exo:g10:numbers:1}
$\frac{15}{3} = 5 \in \N$. \quad $-7 \in \Z$. \quad
$\frac{22}{7} \in \Q$ (bukan \omterm{def:g10:numbers:sets}{bilangan bulat}: $22 = 7 \times 3 + 1$).
\quad $\sqrt 9 = 3 \in \N$. \quad $\sqrt{10} \in \R$ (\omterm{ex:g10:numbers:classify}{irasional}, sebab
$10$ bukan kuadrat sebuah \omterm{def:g10:numbers:sets}{bilangan rasional} --- diterima di sini, sejalan
dengan \cref{thm:g10:numbers:sqrt2}). \quad $-2.4 = -\frac{24}{10} \in \Q$.
\quad $0 \in \N$.
\end{solution}

\begin{solution}{exo:g10:numbers:2}
(a) $\intco{-2}{5}$: titik penuh di $-2$, titik kosong di $5$.
(b) $\intoo{3}{+\infty}$: titik kosong di $3$, arsiran ke kanan.
(c) $\intoc{-\infty}{-1}$: arsiran dari kiri sampai titik penuh di
$-1$.
(d) ``jarak dari $x$ ke $2$ paling banyak $3$'' berarti $\abs{x - 2} \leq 3$,
yaitu $x \in \intcc{-1}{5}$: titik penuh di $-1$ dan $5$.
\end{solution}

\begin{solution}{exo:g10:numbers:3}
(a) Kedua \omterm{def:g10:numbers:interval}{interval} bertumpang tindih antara $0$ dan $2$:
$I \cap J = \intcc{0}{2}$ ($0 \in J$ dan $2 \in I$, dan keduanya juga
anggota himpunan yang lain), lalu $I \cup J = \intco{-3}{4}$.

(b) $I \cap J$ adalah himpunan $x$ dengan $-2 \leq x$ dan $x < 1$:
$\intco{-2}{1}$. \omterm{def:g10:numbers:interunion}{Gabungannya} meliputi segalanya: $I \cup J = \R$.
\end{solution}

\begin{solution}{exo:g10:numbers:4}
$\abs{-6} = 6$; \quad $\abs{4 - 9} = \abs{-5} = 5$; \quad
$\abs{-3 - 5} = \abs{-8} = 8$; \quad
$\sqrt 2 > 1$ sehingga $\abs{\sqrt2 - 1} = \sqrt2 - 1$; \quad
$1 - \sqrt2 < 0$ sehingga $\abs{1 - \sqrt2} = \sqrt2 - 1$ juga (sebuah
bilangan dan lawannya mempunyai \omterm{def:g10:numbers:abs}{nilai mutlak} yang sama).
\end{solution}

\begin{solution}{exo:g10:numbers:5}
$\abs{x} = 5$: jarak $5$ dari $0$, jadi $x = 5$ atau $x = -5$.

$\abs{x - 1} = 3$: jarak $3$ dari $1$, jadi $x = 4$ atau $x = -2$.

$\abs{x - 4} \leq 1$: jarak paling banyak $1$ dari $4$, jadi
$x \in \intcc{3}{5}$.

$\abs{x + 2} < 3$: jarak tegas kurang dari $3$ dari $-2$, jadi
$x \in \intoo{-5}{1}$.
\end{solution}

\begin{solution}{exo:g10:numbers:6}
Setiap \omterm{def:g10:numbers:interval}{interval} dinyatakan oleh pusatnya $a$ (titik tengah kedua batas) dan
jari-jarinya $r$ (setengah panjangnya).

$\intcc{2}{8}$: $a = 5$, $r = 3$, jadi $\abs{x - 5} \leq 3$.

$\intoo{-5}{1}$: $a = -2$, $r = 3$, batas terbuka, jadi $\abs{x + 2} < 3$.

$\intcc{-7}{-3}$: $a = -5$, $r = 2$, jadi $\abs{x + 5} \leq 2$.
\end{solution}

\begin{solution}{exo:g10:numbers:7}
Misalkan $x = 0.272727\dots$ Maka $100x = 27.2727\dots$, lalu dikurangkan:
\[
100x - x = 27.2727\dots - 0.2727\dots = 27,
\]
sehingga $99x = 27$ dan $x = \frac{27}{99} = \frac{3}{11}$, yaitu hasil bagi
dua \omterm{def:g10:numbers:sets}{bilangan bulat}: $x$ bersifat rasional.
\end{solution}

\begin{solution}{exo:g10:numbers:8}
\emph{1. Benar:} menjumlahkan \omterm{def:g10:numbers:sets}{bilangan bulat} (bilangan utuh positif atau
negatif) selalu menghasilkan \omterm{def:g10:numbers:sets}{bilangan bulat}.

\emph{2. Salah:} $\frac{1}{2}$ adalah hasil bagi \omterm{def:g10:numbers:sets}{bilangan bulat} $1$ dan $2$
dan bukan \omterm{def:g10:numbers:sets}{bilangan bulat}.

\emph{3. Benar:} $\frac pq + \frac{p'}{q'} = \frac{pq' + p'q}{qq'}$ lagi-lagi
merupakan hasil bagi dua \omterm{def:g10:numbers:sets}{bilangan bulat} (dengan penyebut taknol).

\emph{4. Benar:} andaikan $r$ rasional, $t$ \omterm{ex:g10:numbers:classify}{irasional}, dan $r + t = s$
juga rasional. Maka $t = s - r$ menjadi selisih dua \omterm{def:g10:numbers:sets}{bilangan rasional},
sehingga rasional (menurut 3, dengan $-r$) --- kontradiksi. Jadi $r + t$
bersifat \omterm{ex:g10:numbers:classify}{irasional}.
\end{solution}

\begin{solution}{exo:g10:numbers:9}
Mengalikan $1.414 \leq \sqrt2 \leq 1.415$ dengan $2 > 0$:
$2.828 \leq 2\sqrt2 \leq 2.830$.

Menambahkan $3$: $4.414 \leq \sqrt2 + 3 \leq 4.415$.

Mengalikan dengan $-1 < 0$ membalik arah pertidaksamaan:
$-1.415 \leq -\sqrt2 \leq -1.414$.
\end{solution}

\begin{solution}{exo:g10:numbers:10}
Pertama fakta bantunya. Bagilah $p$ oleh $3$: sisanya $0$, $1$ atau
$2$, yaitu $p = 3k$, $p = 3k+1$ atau $p = 3k+2$. Dikuadratkan:
\[
(3k)^2 = 3(3k^2), \quad
(3k+1)^2 = 3(3k^2 + 2k) + 1, \quad
(3k+2)^2 = 3(3k^2 + 4k + 1) + 1 .
\]
Hanya yang pertama merupakan kelipatan $3$: jika $p^2$ kelipatan $3$, maka
$p$ juga demikian.

Kini andaikan $\sqrt3 = \frac pq$ sudah paling sederhana. Mengkuadratkan
memberi $p^2 = 3q^2$, sehingga $p^2$ kelipatan $3$, sehingga $p = 3k$. Lalu
$9k^2 = 3q^2$, sehingga $q^2 = 3k^2$ dan $q$ pun kelipatan $3$ --- padahal
pecahan $\frac pq$ tadi sudah paling sederhana: kontradiksi. Jadi
$\sqrt3$ bersifat \omterm{ex:g10:numbers:classify}{irasional}.
\end{solution}

\begin{solution}{pb:g10:numbers:1}
\textbf{1.} $-7 \in \Z$; \ $\frac{13}{4} \in \Q$; \
$\sqrt{16} = 4 \in \N$; \ $0.121212\ldots = \frac{12}{99} =
\frac{4}{33} \in \Q$ (soal akhir pekan tentang desimal berulang di jilid
sebelumnya); \ $\sqrt8 = 2\sqrt2$ bersifat \omterm{ex:g10:numbers:classify}{irasional}, jadi himpunan
terkecilnya $\R$; \ $\pi$: $\R$.

\textbf{2.} $\frac pq + \frac rs = \frac{ps + qr}{qs}$ dan
$\frac pq \times \frac rs = \frac{pr}{qs}$: keduanya hasil bagi
\omterm{def:g10:numbers:sets}{bilangan bulat} dengan penyebut taknol --- jadi rasional.

\textbf{3.} (a) Jika $r$ rasional, $x$ \omterm{ex:g10:numbers:classify}{irasional} dan
$r + x = q$ juga rasional, maka $x = q - r$ menjadi
selisih dua \omterm{def:g10:numbers:sets}{bilangan rasional}, sehingga rasional (pertanyaan 2):
kontradiksi. (b) Jika $r \neq 0$ dan $r x = q$ rasional, maka
$x = \frac qr$ rasional: kontradiksi.

\textbf{4.} $\sqrt2$ dan $-\sqrt2$ \omterm{ex:g10:numbers:classify}{irasional} dengan jumlah $0$;
$\sqrt2$ dan $\sqrt2$ berhasil kali $2$. Hasil rasional dari
bahan \omterm{ex:g10:numbers:classify}{irasional}: \omterm{ex:g10:numbers:classify}{bilangan irasional} tidak membentuk kerajaan
yang tertutup.

\textbf{5.} $\sqrt8 = 2\sqrt2$, sehingga
$\sqrt2 + \sqrt8 = 3\sqrt2$: \omterm{ex:g10:numbers:classify}{irasional} (pertanyaan 3b, dengan pengali
$3$). Dan $\sqrt2 \times \sqrt8 = \sqrt{16} = 4 \in \N$: hasil
kali dua \omterm{ex:g10:numbers:classify}{bilangan irasional} mendarat di kerajaan yang paling
kecil.

\textbf{6.} $3.475$, lalu $3.471$ (atau $3.4701$, $3.47001$,
\dots): menambahkan angka desimal menghasilkan \omterm{def:g10:numbers:sets}{bilangan rasional}
baru di antara keduanya tanpa henti.

\textbf{7.} Kelipatan $10^{-n}$ berbaris di sepanjang garis dengan
langkah $10^{-n} < g$. Kelipatan pertama yang tegas lebih besar
daripada $a$ --- kelipatan itu ada, sebab kelipatan tersebut akhirnya
melampaui $a$ --- terletak paling jauh satu langkah setelah $a$, jadi
sebelum $a + g = b$: tegas di antara $a$ dan $b$. Titik kisi adalah
bilangan desimal, jadi rasional: setiap \omterm{def:g10:numbers:interval}{interval} yang panjangnya positif
memuat satu titik semacam itu.

\textbf{8.} Di antara $a$ dan \omterm{def:g10:numbers:sets}{bilangan rasional} $r_1$ pada
pertanyaan 7 terdapat (pertanyaan 7 lagi) \omterm{def:g10:numbers:sets}{bilangan rasional} $r_2$; di antara
$a$ dan $r_2$ terdapat \omterm{def:g10:numbers:sets}{bilangan rasional} $r_3$; dan seterusnya: tak hingga
banyak, semuanya berbeda, semuanya di dalam \omterm{def:g10:numbers:interval}{interval} semula.

\textbf{9.} $r + \frac{\sqrt2}{10^k}$ adalah jumlah sebuah \omterm{def:g10:numbers:sets}{bilangan
rasional} dan sebuah \omterm{ex:g10:numbers:classify}{bilangan irasional} ($\frac{\sqrt2}{10^k}$ \omterm{ex:g10:numbers:classify}{irasional}
menurut pertanyaan 3b), jadi \omterm{ex:g10:numbers:classify}{irasional}. Karena
$\frac{\sqrt2}{10^k} < \frac{2}{10^k}$ mengecil di bawah batas mana pun,
untuk $k$ cukup besar $r + \frac{\sqrt2}{10^k}$ masih terletak
sebelum $b$: sebuah \omterm{ex:g10:numbers:classify}{bilangan irasional} di dalam \omterm{def:g10:numbers:interval}{interval} --- dan dengan
mengubah $k$ diperoleh tak hingga banyak.

\textbf{10.} Tidak ada \omterm{def:g10:numbers:sets}{bilangan real} positif terkecil: jika $s > 0$
adalah bilangan itu, \omterm{def:g10:numbers:interval}{interval} $(0, s)$ masih memuat \omterm{def:g10:numbers:sets}{bilangan rasional}
(pertanyaan 7) yang positif dan lebih kecil daripada $s$. Tidak ada
bilangan tepat setelah $3$: setiap calon $c > 3$ menyisakan \omterm{def:g10:numbers:interval}{interval}
$(3, c)$ yang tidak kosong --- permainan pada soal akhir pekan tentang
tiadanya bilangan berikutnya di jilid sebelumnya, kini menjadi teorema
tentang $\R$.

\textbf{11.} $\abs{x - 5} = 2$: $x = 3$ atau $x = 7$.
$\abs{x - 5} < 2$: $x \in \intoo{3}{7}$.
$\abs{x + 1} \geq 3$: jarak ke $-1$ paling sedikit $3$:
$x \in \intoc{-\infty}{-4} \cup \intco{2}{+\infty}$.

\textbf{12.} $\abs{d - 80} \leq 0.05$, yaitu
$d \in \intcc{79.95}{80.05}$. Piston berukuran $79.97$ lolos;
yang berukuran $80.06$ gagal, meleset seperseratus milimeter.

\textbf{13.} $(3, -5)$: $\abs{-2} = 2 \leq 8$. $(-2, -7)$:
$\abs{-9} = 9 = 2 + 7$: kesamaan. Bertanda sama: $\abs{a + b}$ adalah
jumlah kedua jarak, sama dengan $\abs a + \abs b$. Bertanda berlawanan:
langkahnya berbalik arah, $\abs{a + b}$ menjadi
\emph{selisih} kedua jarak, tegas lebih kecil daripada
jumlahnya (kecuali salah satunya $0$). Kesamaan berlaku tepat ketika $a$
dan $b$ bertanda sama atau salah satunya nol.

\textbf{14.} Penyelesaiannya adalah titik yang berjarak sama dari $2$
dan $-4$: titik tengahnya, $x = -1$. Titik itu adalah garis sumbu pada
soal akhir pekan tentang dua cermin di jilid sebelumnya, yang mengerut
ke satu dimensi: sebuah titik tunggal.

\textbf{15.} $\sqrt{x^2} = \abs{x}$, bukan $x$: untuk $x = -3$,
$\sqrt{9} = 3 = \abs{-3} \neq -3$. Lalu $x^2 < 9$ terbaca
$\sqrt{x^2} < 3$, yaitu $\abs x < 3$:
$x \in \intoo{-3}{3}$.

\textbf{16.} Pengukuran itu hanya menegaskan bahwa panjangnya terletak
di dalam \omterm{def:g10:numbers:interval}{interval} $\intcc{1.41420}{1.41422}$ --- dan menurut pertanyaan
7 dan 9 \emph{\omterm{def:g10:numbers:interval}{interval} itu memuat tak hingga banyak \omterm{def:g10:numbers:sets}{bilangan rasional}
dan tak hingga banyak \omterm{ex:g10:numbers:classify}{bilangan irasional}}. Hal yang sama berlaku bagi
setiap toleransi di kemudian hari, sekecil apa pun. Tidak ada pengukuran
yang mampu memisahkan kedua keluarga itu; hanya bukti tentang panjang
yang eksak --- seperti bukti untuk diagonal persegi satuan (soal akhir
pekan tentang sifat \omterm{ex:g10:numbers:classify}{irasional} di jilid sebelumnya) --- yang mampu.

\textbf{17.} \omterm{def:g10:numbers:approx}{Hampiran} itu mendekati $\sqrt2$ dengan galat di bawah
$10^{-1}, 10^{-2}, 10^{-3}, \dots$: \omterm{def:g10:numbers:sets}{bilangan rasional} mendesak
$\sqrt2$ pada setiap skala. Fakta umumnya: setiap \omterm{def:g10:numbers:sets}{bilangan real} dapat
didekati sedekat yang diinginkan oleh \omterm{def:g10:numbers:sets}{bilangan rasional} (pemenggalan
desimalnya) --- kepadatan sekali lagi, dilihat dari sisi sasarannya.

\textbf{18.} $a + b \in \intcc{2.4 + 1.1}{2.6 + 1.3} =
\intcc{3.5}{3.9}$: ketakpastian $\pm 0.2$ (galat bertambah).
$a \times b \in \intcc{2.4 \times 1.1}{2.6 \times 1.3} =
\intcc{2.64}{3.38}$: ketakpastian sekitar $\pm 0.37$ di sekitar
$3$ --- pada perkalian yang bertambah adalah galat relatifnya, jadi
ketelitian merosot lebih cepat.

\textbf{19.} Tulislah nilai sebenarnya sebagai $a = 2.5 + e_1$,
$b = 1.2 + e_2$ dengan $\abs{e_1}, \abs{e_2} \leq 0.1$. Maka
galat jumlahnya adalah $\abs{e_1 + e_2} \leq \abs{e_1} +
\abs{e_2} \leq 0.2$: aturan insinyur itu \emph{adalah} ketaksamaan
segitiga pada pertanyaan 13.

\textbf{20.} Garis \omterm{def:g10:numbers:sets}{bilangan real} tidak berlubang dan tidak bertetangga:
setelah suatu bilangan, tidak ada bilangan ``berikutnya'' (pertanyaan
10). Titik-titiknya terbelah menjadi \omterm{def:g10:numbers:sets}{bilangan rasional} --- tepatnya
desimal yang berhenti atau berulang (soal akhir pekan tentang desimal
berulang di jilid sebelumnya) --- dan \omterm{ex:g10:numbers:classify}{bilangan irasional}, dan kedua
keluarga itu terjalin begitu rapat sehingga setiap \omterm{def:g10:numbers:interval}{interval} memuat tak
hingga banyak dari masing-masing (pertanyaan 8 dan 9); itulah sebabnya
tidak ada pengukuran, melainkan hanya bukti, yang dapat membedakannya
(pertanyaan 16). Dan kejutan terakhir, yang ditinggalkan sebagai janji:
\omterm{ex:g10:numbers:classify}{bilangan irasional} lebih banyak daripada \omterm{def:g10:numbers:sets}{bilangan rasional} --- bukan
dengan membilang satu per satu, melainkan dalam pengertian yang tepat
tentang membandingkan ketakhinggaan, sebuah teori yang dibangun pada
jilid universitas.

\end{solution}
